\documentclass[parskip=full]{scrreprt}

\RedeclareSectionCommand[pagestyle=plain,afterskip=10pt plus 1pt]{chapter}
\setkomafont{chapter}{\mdseries\headingfont\fontsize{40}{40}\selectfont\color{black!80}}
\setkomafont{pageheadfoot}{\normalsize}

\def\pnumbox#1{#1\hspace*{9cm}}
\DeclareTOCStyleEntry[
  indent=0pt,
  entrynumberformat=\textcolor{oldred},
  numwidth=2em,
  linefill=\hfill,
  pagenumberbox=\pnumbox,
  pagenumberformat=\itshape
]{tocline}{section}

\usepackage[english]{babel}

\usepackage{fontspec}

\setmainfont{Source Sans Pro}[
  ItalicFont = Source Sans Pro Italic,
  BoldFont = Source Sans Pro Bold,
  BoldItalicFont = Source Sans Pro Bold Italic,
  FontFace = {lt}{n}{Source Sans Pro Light},
  FontFace = {lt}{it}{Source Sans Pro Light Italic},
  FontFace = {sb}{n}{Source Sans Pro Semibold},
  FontFace = {sb}{it}{Source Sans Pro Semibold Italic},
  Numbers = Proportional,
  Ligatures = Common
]
\DeclareRobustCommand{\ltseries}{\fontseries{lt}\selectfont}
\DeclareRobustCommand{\sbseries}{\fontseries{sb}\selectfont}
\DeclareTextFontCommand{\textlt}{\ltseries}
\DeclareTextFontCommand{\textsb}{\sbseries}

\newfontfamily\headingfont{Fredericka the Great}

\usepackage[pass]{geometry}
\newgeometry{twoside,inner=20mm,outer=40mm,top=20mm,bottom=40mm}

\usepackage{url}
\urlstyle{same}

\usepackage{microtype}
\microtypesetup{verbose=silent}

\usepackage{booktabs,array,longtable}
\newcolumntype{L}[1]{>{\raggedright\let\newline\\\arraybackslash\hspace{0pt}}p{#1}}

\usepackage{graphicx}

\usepackage{xcolor}
\definecolor{oldred}{rgb}{.8313,0,0}

\usepackage{pdfpages}

\usepackage{multicol}

\usepackage{scrlayer-scrpage}
\pagestyle{scrheadings}
\clearscrheadfoot
\cfoot[\thepage]{\thepage}
\pagenumbering{roman}


\makeatletter

\newcommand\fancytitlehead{
	\headingfont%
	\fontsize{80}{80}\selectfont\textcolor{black!80}{\@lastname.}\\[15pt]%
	\fontsize{60}{60}\selectfont\@ifundefined{@shorttitle}{\@title}{\@shorttitle}.%
}


\def\firstname#1{\def\@firstname{#1}}
\def\lastname#1{\def\@lastname{#1}}
\def\shorttitle#1{\def\@shorttitle{#1}}
\def\namesuffix#1{\def\@namesuffix{#1}}
\def\instrumentation#1{\def\@instrumentation{#1}}
\def\parts#1{\def\@parts{#1}}

\firstname{\relax}
\lastname{\relax}
\namesuffix{\relax}
\instrumentation{\relax}
\parts{\relax}


\def\maketitle{%
\begin{titlepage}%
	\Large%
	{\@titlehead}%
	\vfill%
	{\strut\@firstname}\\%
	{\sbseries\color{oldred}\strut\@lastname}\\%
	{\strut\@namesuffix}%
	\vfill%
	{\sbseries\@title}\\%
	{\@subtitle}\\[\baselineskip]%
	{\itshape\@instrumentation}%
	\vfill%
	{\itshape\@parts}\hspace*{\fill}\raisebox{0pt}[0pt][0pt]{\includegraphics{ees_logo}}%
\end{titlepage}%
}


\newif\iftemplate\templatetrue
\newif\ifprintreport\printreportfalse
\directlua{
scores = {
  cl1 = "Clarinetto I",
  cl2 = "Clarinetto II",
  fag1 = "Fagotto I",
  fag2 = "Fagotto II",
  ottoni = "Corno I, II\string\\newline Tromba I, II\string\\newline Timpani",
  trb123 = "Trombone I, II, III",
  vl1 = "Violino I",
  vl2 = "Violino II",
  vla = "Viola",
  coro = "Coro",
  org = "Organo",
  b = "Bassi",
  full_score = "Full Score"
}

last_arg = arg[\string#arg]
texio.write("Last argument: " .. last_arg)
if not (scores[last_arg] == nil) then
  tex.print("\string\\def\string\\lypdfname{" .. last_arg .. "}")
  tex.print("\string\\parts{" .. scores[last_arg] .. "}")
  if (last_arg == "full_score") then
    tex.print("\string\\printreporttrue")
  end
end
}

\@ifundefined{lypdfname}{%
  \templatefalse
  \printreporttrue
  \parts{Draft}
}{\templatetrue}

\def\@firstoffour#1#2#3#4{#1}
\def\@secondoffour#1#2#3#4{#2}
\def\@thirdoffour#1#2#3#4{#3}
\def\@lastoffour#1#2#3#4{#4}
\def\lyrefnumber#1{\expandafter\@setref\csname r@#1\endcsname\@firstoffour{#1}}
\def\lyreftitle#1{\expandafter\@setref\csname r@#1\endcsname\@secondoffour{#1}}
\def\lyreftitledansk#1{\expandafter\@setref\csname r@#1\endcsname\@thirdoffour{#1}}
\def\lyrefpage#1{\expandafter\@setref\csname r@#1\endcsname\@lastoffour{#1}}
\def\lyref#1{%
  \begingroup%
  \@tempdima.5\linewidth\relax%
  \advance\@tempdima.5\columnsep\relax%
  \sbseries%
  \makebox[0pt][r]{\color{oldred}\lyrefnumber{#1}\hspace*{1.5em}}%
  \makebox[\@tempdima][l]{\lyreftitle{#1}\Dotfill}%
  \textit{\lyreftitledansk{#1}}%
  \Dotfill\lyrefpage{#1}%
  \endgroup%
}

\def\Dotfill{\leavevmode\cleaders\hb@xt@ .75em{\hss .\hss }\hfill \kern \z@}

\makeatother

\input{../lilypond.ref}




\begin{document}

\titlehead{\fancytitlehead}
\firstname{Friedrich Ludwig Æmilius} % Libretto Jens Immanuel Baggensen
\lastname{Kunzen}
\title{Skabningens Halleluja}
\shorttitle{Skabningens\\Halleluja}
\subtitle{Das Halleluja der Schöpfung\\(DK-Kk mu 6506.1133)}
\instrumentation{2 S, A, T, B (coro), 2 fl, 2 ob, 2 cl, 2 fag, 2 cor, 2 tr, 3 trb, timp,\\vl solo, 2 vl, vla, 2 vlc, b}
\maketitle


\thispagestyle{empty}

\vspace*{\fill}

\raisebox{-4mm}{\includegraphics{byncsaeu}}\hspace*{1em}Wolfgang Esser-Skala, 2020

© 2020 by Wolfgang Esser-Skala. This edition is licensed under the Creative Commons Attribution-NonCommercial-ShareAlike 4.0 International License. To view a copy of this license, visit \url{http://creativecommons.org/licenses/by-nc-sa/4.0/}. 

Music engraving by LilyPond 2.18.0 (\url{http://www.lilypond.org}).\\
Front matter typeset with Source Sans Pro and Fredericka the Great.

\textit{First version, April 2020}

\vspace*{2cm}

\ifprintreport
\chapter*{Critical Report.}

This edition bases upon a manuscript in the Det Kongelige Bibliotek på Slotsholmen – Den Sorte Diamant. The digital version of the manuscript is available at \url{http://img.kb.dk/ma/kunzen/SkaHal_m.pdf} (siglum mu 6506.1133).

In general, this edition closely follows the manuscript. Any changes that were introduced by the editor are indicated by italic type (dynamics and directions), parentheses (expressive marks) or dashes (slurs and ties). Accidentals are used according to modern conventions. Asterisks denote changes that are clarified in the detailed remarks below.

\bigskip

\footnotetext[1]{Abbreviations: A, alto; b, violoncello and bass; B, basso; cl, clarinet; cor, horn; fag, bassoon; ob, oboe; Ms, manuscript; r, rest; S, soprano; T, tenore; timp, timpani; tr, trumpet; trb, trombone; vl, violin; vla, viola.}

\begin{longtable}{lll L{10cm}}
	\toprule
	\itshape Mov. & \itshape Bar & \itshape Staff & \itshape Note \\
	\midrule \endhead
	–  & –   & org   & All bass figures were supplemented by the editor. \\
	1  & 58  & timp  & 2nd half of bar in Ms: c4–c4 \\
	   & 31  & B     & 4th quarter: grace note missing in Ms \\
	4  & 43  & cor 2 & 4th quarter in Ms: c″4 \\
	5  & 16  & fag 2 & grace note missing in Ms \\
	   & 34  & T     & grace note missing in Ms \\
	   & 44  & vlc   & bar in Ms: G8–d8–f8–e4. \\
	   & 108 & fag 1 & 3rd quarter in Ms: r8–d8 \\
	   & 149 & fag 1 & bar in Ms: r1 \\
	\bottomrule
\end{longtable}


This edition has been compiled and checked with utmost diligence. Nevertheless, errors and mistakes cannot be totally excluded. Please report any errors and mistakes to \url{wolfgang@esser-skala.at} or create an issue or pull request on the edition’s GitHub page (https://github.com/skafdasschaf/kunzen-skabningens-halleluja). Your help will be greatly appreciated.

\bigskip
\textit{Salzburg, April 2020\\
Wolfgang Esser-Skala}




\chapter*{Contents. Lyrics.}


\begin{multicols}{2}[\lyref{brichnatur}]
Brich, Natur, in Loblied aus!\\
Lobt Jehoven, alle Zonen,\\
lobet, alle Nationen,\\
lobet Gott!\\
Preiset ihn in vollen Chören,\\[\baselineskip]
Himmel, Erd und Meer!\\
Die ihr athmet, preist den Guten!\\
Ihn, der schuf des Lichtes Gluthen,\\
Ihn, der Odem gab!\\
Hoch in frohe Jubeltöne\\
schweb empor, Gesang, zu Gott!\\
Der gesammten Schöpfung Söhne\\
brechet aus in Jubeltöne,\\
lobet Gott!\\
Der gesammten Schöpfung Söhne\\
singt im heiligsten der Töne,\\
lobet Gott!\columnbreak

\itshape
Bryd, o Støv, i Lovsang ud!\\
Lover Herren alle Zoner!\\
Lover alle Nationer!\\
Lover Gud!\\
Priiser ham, den Eene Store!\\
Priiser Ham i fulde Chore,\\
Himmel, Jord, og Hav!\\
Alt hvad aander, nynne, stamme\\
Ham, som tændte Lysets Flamme!\\
Ham, som Aande gav!\\
Høit i glade Jubeltoner\\
bryde Livets Lovsang [\textit{or}: Vellyst] ud!\\
Alt det Skabtes [\textit{or}: Alle Kloders] Millioner,\\
toner i foreente Toner,\\
toner Gud!
\end{multicols}

\begin{multicols}{2}[\lyref{dieoede}]
Die Öde starrt,\\
von keinem Strahl erhellt,\\
im Schlund der Nacht war Leere, Dunkel, Tod.\\
Da schallt durchs Chaos schaffend dein Gebot:\\
„Sey Welt!“\\
Ein Sonnenheer entsprang aus alter Nacht\\
in Licht hervor.\\
Der Neugeschaffnen Hymne tönt empor:\\
„Halleluja wir leben!\\
Halleluja! du bist! du warst! und bleibest ewig\\
Herr, unser Gott!“\columnbreak

\itshape
Alt hylled laae [\textit{or}: slumred dybt]\\
i Muelighedens Skiød,\\
i Nattens Svælg var Tomhed, Mørke, Død.\\
Men giennem Evighedens Chaos lød\\
dit Bliv!\\
Og Verdener sprang frem af Intets Skiød\\
til Lys og Liv.\\
Nyskabte Væsners Studsen rundt brød ud:\\
„Halleluja! vi ere!\\
Halleluja! du er, du var, og du skal være\\
vor Gud!”
\end{multicols}
\clearpage

\begin{multicols}{2}[\lyref{vomschlaf}]
Vom Schlaf im Schattenwald erwacht,\\
sieht, wenn der Morgen wieder lacht,\\
was lebt, sich deine Herrlichkeit erneun!\\
Sieh! Vögel froh auf leichten Schwingen\\
hinauf in goldne Wolken dringen,\\
ihr Loblied wiederhallt in Flur und Hain!\\
Weit strahlet wachsend deine lichte\\
\hspace*{1em}Morgenröthe!\\
Mild glänzt der Thau, die Rose glüht,\\
sanft weht die Luft, der Bach entflieht,\\
und Haingesang ertönt zur Schöpfungs Feier,\\
milder Thau, Rosenduft,\\
Haingesang, kühle Luft\\
begrüßt den Tag, erweckt zur Morgenfeier.\columnbreak

\itshape
Fra Skovens [\textit{or}: Lundens] skyggefulde Lye\\
seer vækket Liv, hver Morgen nye,\\
i Mark [\textit{or}: Skov] og Eng din Herlighed fremgrye!\\
See! Fuglen fro, paa lette Vinger\\
sig op i Luftens Purpur svinger,\\
Din [\textit{or}: Dens] Lovsang dirrer høit i gyldne Skye!\\
Vidt straaler voxende din lyse\\
\hspace*{1em}Morgenrøde!\\
Og Duggens Blik, og Blomstrets Smiil,\\
og Vindens Vift, og Bækkens Iil,\\
i yndig Kappen stræber den imøde,\\
Duggens Blik, Blomstrets Smiil,\\
Vindens Vift, Bækkens Iil\\
i yndig Kappen stræber den imøde.
\end{multicols}

\begin{multicols}{2}[\lyref{durollest}]
Du rollest auf der Dämmrung Flor!\\
Die Thräne trocknest du,\\
die thauend Blumen tränkte!\\
Dein Himmel graut nicht, wie zuvor –\\
des Tages lichter Held tritt auf in vollem Glanze!\\
Umarmend die blühende Braut,\\
entzückt von des Lichtes belebender Wonne,\\
erzählet die Schöpfung es laut:\\
Das ist Gott!\\
Halleluja! wir leben!\\
Halleluja! du bist, du warst, und bleibest ewig,\\
Herr, unser Gott!\columnbreak

\itshape
Du ruller bort Naturens Slør!\\
Du tørrer Duggens Graad,\\
som Jordens Ansigt væder!\\
Din Himmel skiules ei, som før —\\
see! Dagens lyse Helt i fulde Glands fremtræder!\\
Han favner sin rødmende Brud,\\
og tryllet af Lysets oplivende Glæder\\
den hele Natur bryder ud:\\
Det er Gud!\\
Halleluja! vi ere!\\
Halleluja! du er, du var, og du skal være,\\
vor Gud!
\end{multicols}
\newpage
\begin{multicols}{2}[\lyref{ichhoerte}]
Ich hörte Haingesang,\\
die Rose that sich auf, das Laub erbebte;\\
ich sah das Morgenroth\\
und fühlte Lust, die neu mein Herz belebte.

Doch einsam war ich, meine Brust\\
fand leer die nicht getheilte Lust.\\
Ich fand die junge Ros im Morgenlicht;\\
ich lächelt ihr; doch sie verstand mich nicht.

Da hörtest Schöpfer du der Sehnsucht Regen;\\
und sieh!\\
Holdlächelnd kam die Gattin mir entgegen!\\
Nun blickt ich froh ins Morgenroth,\\
und rufte: Gut ist Gott!

Ich sah der Sterne Heer\\
im festlichen Triumph den Tag begleiten;\\
ich sah der Sonne Sieg,\\
mein Geist verlohr sich in des Himmels\\
\hspace*{1em}Weiten.

Doch einsam war ich, meine Brust\\
fand Wehmut nur in halber Lust.\\
Ich rief den Bäumen, sich mit mir zu freun;\\
doch mich verstand kein Baum im ganzen Hain.

Da hörtest du der Thräne leises Flehen;\\
und sieh!\\
Mein ander Ich trat her von jenen Höhen –\\
O! Gut ist Gott! rief er mir zu,\\
ich rufte: Gott ist gut!

Ich sah des Morgens Purpurzelt –\\
ich sahs, er kam, des Tages Held –\\
doch was ist, Morgenroth, dein Blick?\\
Wie weicht der Sonne Gluth zurück?\\
Was bist du schönes Thal –\\
was bist du hoher Himmel –\\
vorm Feuerblick des Manns?\\
Ach, vor der Gattin Huld?

O, was ist aller Herrlichkeiten Fülle\\
vor dir, der Liebe Gluth?\\
Vor unsrer Herzen Lust?

Preis, Liebe, dir, wir leben!\\
Preis, Liebe, dir! du bist, du warst\\
und bleibest ewig,\\
o Liebe, unsre Lust!

\columnbreak\itshape
Jeg hørte Lærken slaae,\\
og Rosen aabned sig, og Løvet bæved;\\
jeg Morgenrøden saae,\\
og, vakt til Lyst, mit Hierte høit sig hæved.

Men eene var jeg, og mit Bryst\\
fandt Savn i sin udeelte Lyst.\\
Jeg Roser fandt paa min bestraalte Vei;\\
de rødmede; men de forstod mig ei.

Da saae du, Skaber! mine Taarer rinde;\\
og see!\\
De sendte mig den smilende Mandinde!\\
Jeg sank i Fryd ved Eegens Rod,\\
og raabte: Gud er god!

Jeg Lysets Hære saae\\
med Dagen i Triumph til Jorden stimle;\\
jeg Solen seire saae,\\
og Tanken tabte sig i Himles Himle.

Men eene var jeg, og mit Bryst\\
fandt Veemod i sin halve Lyst.\\
Jeg iilte længselfuld hvert Træ imod;\\
men intet Træ i Lunden mig forstod.

Da saae du, Fader! Taaren i mit Øye;\\
og see!\\
Min anden Jeg kom frem bag Skovens Høye —\\
O! Gud er god! saa brød han ud,\\
jeg svarte: God er Gud!

Jeg Østens lyse Purpur saae —\\
jeg skued Dagens Helt fremgaae —\\
men hvad er Morgenrødens Smiil?\\
Men hvad er Solens Straalers Ild?\\
Hvad er du, skiønne Jord —\\
hvad er du, høie Himmel —\\
mod Mandens Tankeblik?\\
Imod Mandindens Smiil?

O hvad er alle Herligheders Vrimmel\\
mod Flammen i vort Bryst?\\
Mod vore Hierters Lyst?

O Kierlighed! vi ere!\\
O Kierlighed! du er, du var,\\
og du skal være\\
vor Lyst!
\end{multicols}



\cleardoublepage
\fi

\iftemplate
\includepdf[pages=-]{../out/\lypdfname.pdf}
\fi



\end{document}